% RACER-GT Python API 參考（繁體中文）。
\documentclass[12pt,a4paper,oneside]{article}
% Traditional Chinese setup. Compile with XeLaTeX.
\usepackage{fontspec}
\usepackage{xeCJK}
\setmainfont{Noto Serif}
\setsansfont{Noto Sans}
\setmonofont{Noto Sans Mono}
\setCJKmainfont{Noto Serif CJK TC}
\setCJKsansfont{Noto Sans CJK TC}
\setCJKmonofont{Noto Sans CJK TC}
\XeTeXlinebreaklocale "zh"
\XeTeXlinebreakskip = 0pt plus 1pt

% Single source of truth for version and date across every RACER-GT document.
% Regenerate nothing by hand: bump here, rebuild, and every PDF agrees.
\newcommand{\racerversion}{1.1.0}
\newcommand{\racerdate}{2026-07-27}
\newcommand{\racerauthor}{Yi-Hao Lai}
\newcommand{\raceraffil}{Department of Finance, Da-Yeh University}
\newcommand{\raceremail}{yhlai@mail.dyu.edu.tw}
\newcommand{\racerrepo}{https://github.com/yhlai0911/RACER-GT}

% Shared package set and mathematical macros for every RACER-GT document,
% in both languages. Language-specific font and theorem-name setup lives in
% _preamble-zh.tex and _preamble-en.tex.
%
% Font policy: the build depends only on the five Noto families below. They are
% installable on Linux CI (fonts-noto-core, fonts-noto-cjk) and on macOS via
% Homebrew casks. Do not introduce a font outside this set without adding it to
% .github/workflows/docs.yml as well, or the PDF build stops being reproducible.

\usepackage[margin=2.35cm,headheight=15pt]{geometry}
\usepackage{amsmath,amssymb,amsthm,mathtools,bm}
\usepackage{booktabs,longtable,tabularx,array,multirow}
\usepackage{graphicx,float,caption,subcaption}
\usepackage{enumitem}
\usepackage{xcolor}
\usepackage{microtype}
\usepackage{setspace}
\usepackage{fancyhdr}
\usepackage{hyperref}
\usepackage[nameinlink,noabbrev]{cleveref}
\usepackage{algorithm}
\usepackage{algpseudocode}
\usepackage{listings}
\usepackage{tikz}
\usetikzlibrary{arrows.meta,positioning,shapes.geometric,fit,calc}

\hypersetup{
  colorlinks=true,
  linkcolor=black,
  citecolor=black,
  urlcolor=blue,
  pdfauthor={\racerauthor}
}

\pagestyle{fancy}
\fancyhf{}
\rhead{v\racerversion}
\cfoot{\thepage}
\setstretch{1.28}
\setlength{\parskip}{0.35em}
\setlist[itemize]{leftmargin=2em,itemsep=0.25em,topsep=0.3em}
\setlist[enumerate]{leftmargin=2.2em,itemsep=0.25em,topsep=0.3em}

% ---------------------------------------------------------------- mathematics
\newcommand{\E}{\mathbb{E}}
\newcommand{\Var}{\operatorname{Var}}
\newcommand{\Cov}{\operatorname{Cov}}
\newcommand{\tr}{\operatorname{tr}}
\newcommand{\diag}{\operatorname{diag}}
\newcommand{\argmin}{\operatorname*{arg\,min}}
\newcommand{\argmax}{\operatorname*{arg\,max}}
\newcommand{\one}{\bm{1}}
\newcommand{\Rset}{\mathbb{R}}
\newcommand{\R}{\mathbb{R}}
\newcommand{\plim}{\operatorname*{plim}}
\newcommand{\RACER}{\textsc{Racer-GT}}
\newcommand{\indic}[1]{\mathbb{I}\!\left\{#1\right\}}
\newcommand{\Sig}{\bm{\Sigma}}
\newcommand{\wvec}{\bm{w}}
\newcommand{\evec}{\bm{e}}
\newcommand{\Zvec}{\bm{Z}}

% Design facets, used identically in both languages so that a reader can move
% between the Chinese and English documents without re-learning notation.
\newcommand{\facetT}{\mathcal{T}}   % historical date
\newcommand{\facetD}{\mathcal{D}}   % collection day
\newcommand{\facetS}{\mathcal{S}}   % collection stream
\newcommand{\facetR}{\mathcal{R}}   % technical replicate

\lstset{
  basicstyle=\ttfamily\footnotesize,
  breaklines=true,
  frame=single,
  rulecolor=\color{gray!40},
  backgroundcolor=\color{gray!5},
  columns=fullflexible,
  keepspaces=true,
  showstringspaces=false
}


\setlength{\parindent}{2em}

\newtheorem{assumption}{假設}[section]
\newtheorem{proposition}{命題}[section]
\newtheorem{corollary}{推論}[section]
\newtheorem{remark}{說明}[section]
\newtheorem{definition}{定義}[section]
\newtheorem{lemma}{引理}[section]

\renewcommand{\proofname}{證明}
\renewcommand{\abstractname}{摘要}
\renewcommand{\contentsname}{目錄}
\renewcommand{\refname}{參考文獻}
\renewcommand{\tablename}{表}
\renewcommand{\figurename}{圖}

\crefname{assumption}{假設}{假設}
\crefname{proposition}{命題}{命題}
\crefname{corollary}{推論}{推論}
\crefname{definition}{定義}{定義}
\crefname{lemma}{引理}{引理}
\crefname{equation}{式}{式}
\crefname{table}{表}{表}
\crefname{figure}{圖}{圖}
\crefname{section}{節}{節}

\lhead{RACER-GT：API 參考}

\title{\bfseries RACER-GT Python API 參考\\[0.4em]
\large 公開工作流程與診斷進入點}
\author{\racerauthor\\\small 大葉大學 財務金融學系}
\date{\racerdate\ \textbullet\ 版本 \racerversion}

\begin{document}
\maketitle

\begin{abstract}
\noindent 頂層 API 是穩定的。較低階的函式仍可用於方法論診斷，但演進速度可能快於 pipeline 介面。凡是「拋錯而非回傳降級結果」之處，本文都會明確指出——那些拒絕是刻意的設計決策，不是缺少錯誤處理。
\end{abstract}

\tableofcontents
\newpage

\section{頂層 API}

\begin{lstlisting}[language=Python]
from racergt import (
    RacerGTConfig, RacerGTPipeline, PipelineResult,
    DesignWeightedResult, fit_design_weighted_consensus,
    generate_collection_schedule, generate_chunk_windows,
)
\end{lstlisting}

\subsection{\texttt{RacerGTConfig}}

巢狀 pydantic 模型，承載所有鎖定設定。所有子模型均使用 \texttt{extra="forbid"}，因此 YAML 中出現無法辨識的鍵會是驗證錯誤，而不是被靜默忽略的錯字。

\begin{lstlisting}[language=Python]
cfg = RacerGTConfig.load_yaml("protocol.yaml")
cfg.save_yaml("protocol.lock.yaml")
digest = cfg.protocol_hash()          # canonical 序列化後的 SHA-256
\end{lstlisting}

\texttt{load\_yaml} 會重算 hash，若檔案中儲存的 \texttt{protocol\_hash} 不符即\textbf{拋錯}。手動編輯 lock 檔因此會明確失敗。要改設定，請改來源設定檔後重新產生。

\subsection{\texttt{generate\_chunk\_windows(start, end, window\_days, step\_days)}}

回傳固定的重疊 chunk manifest：\texttt{chunk\_id}、\texttt{window\_start}、\texttt{window\_end}、\texttt{n\_calendar\_days}。

\subsection{\texttt{generate\_collection\_schedule(config, anchor\_date=None)}}

回傳平衡後的 day $\times$ stream $\times$ replicate 排程，含 stream 順序、規劃時段、決定性的 chunk 順序、固定歷史端點與 protocol hash。

\begin{remark}
輸出中的 \texttt{collection\_day} 是\emph{序數設計層級}（0、1、2……），不是 day offset；\texttt{day\_offset} 與 \texttt{planned\_collection\_date} 是另外的欄位。
\end{remark}

\subsection{\texttt{RacerGTPipeline(config).fit(raw, benchmark=None)}}

依序執行稽核、各 pull 的重疊校準、duplicate 診斷、設計式參考估計量、G-studies、共變異數調整 consensus、（有提供 benchmark 時的）頻率 benchmark、reliability 診斷，以及鎖定的 decision tree。回傳 \texttt{PipelineResult}。

當原始批次未通過 protocol 稽核時拋出 \texttt{ValueError}，除非設定 \texttt{stop\_on\_audit\_error=False}。

\subsection{\texttt{PipelineResult}}

\begin{table}[H]
\centering\small
\begin{tabularx}{\textwidth}{lX}
\toprule
屬性 & 內容\\
\midrule
\texttt{final\_series} & 最終日頻指標、條件式標準誤、95\% 區間\\
\texttt{calibrations} & 各 pull 的重疊圖校準結果\\
\texttt{duplicate\_diagnostics} & exact groups 與 near-duplicate 圖\\
\texttt{gstudies} & \texttt{level}、\texttt{detection}、\texttt{innovation}\\
\texttt{design\_consensus} & 設計式參考估計量\\
\texttt{consensus} & GLS／最小變異 consensus\\
\texttt{benchmark} & 頻率 benchmark 結果，或 \texttt{None}\\
\texttt{reliability} & pairwise、day/stream、收斂診斷\\
\texttt{decision} & PASS／REVIEW／FAIL 與每一條規則\\
\texttt{audit} & protocol 完整性稽核\\
\bottomrule
\end{tabularx}
\end{table}

\texttt{final\_series} 在有 benchmark 時回傳 benchmark 結果，否則回傳 GLS consensus。\texttt{design\_consensus} 永遠不會進入它。

\section{診斷進入點}

\subsection{校準}

\begin{lstlisting}[language=Python]
from racergt.overlap import OverlapGraphCalibrator, huber_location, OverlapGraphError
\end{lstlisting}

\texttt{OverlapGraphCalibrator.fit} 要求恰好一個 \texttt{series\_id} 與一個 \texttt{pull\_id}。當重疊圖不連通且 \texttt{allow\_disconnected} 為否時拋出 \texttt{OverlapGraphError}——因為圖不連通代表相對尺度只被部分識別，此時仍回傳數列等於隱瞞這件事。

\subsection{Duplicates、G-study、consensus}

\begin{lstlisting}[language=Python]
from racergt.duplicates import diagnose_duplicates
from racergt.gstudy import run_gstudy, run_all_gstudies
from racergt.consensus import fit_design_weighted_consensus, fit_gls_consensus
\end{lstlisting}

\texttt{run\_gstudy} 在設計非完全交叉或 replicate 數不等時拋錯，而不是從不平衡設計估出變異成分。\texttt{fit\_gls\_consensus} 要求至少兩個 pull，且每個 pull 的 baseline 平均為正且有限。

\begin{remark}[Duplicate 分量不是等價類]
Near-duplicate 圖的連通分量是連通集合。屬於同一分量不代表其中每一對都超過門檻，因為連通性具遞移性、而 pairwise 規則沒有。
\end{remark}

\subsection{Benchmark 與測量誤差}

\begin{lstlisting}[language=Python]
from racergt.benchmark import temporal_benchmark
from racergt.eiv import reliability_corrected_ols, simex_ols
\end{lstlisting}

\texttt{reliability\_corrected\_ols} 在修正後分母非正時拋錯。該情況代表估計的測量誤差已吃掉全部可觀察的解釋變數變異；此時回傳任何數值都毫無意義。

\subsection{模擬與驗證}

\begin{lstlisting}[language=Python]
from racergt.simulation import SimulationSettings, simulate_racergt_data
from racergt.validation import run_monte_carlo
\end{lstlisting}

\section{檔案 I/O}

\texttt{racergt.io.read\_table} 可讀 CSV、Parquet、XLS/XLSX；\texttt{write\_table} 可寫 CSV、Parquet、XLSX 與 Stata 118 \texttt{.dta}。

\section{命令列}

\begin{table}[H]
\centering\small
\begin{tabularx}{\textwidth}{lXl}
\toprule
指令 & 用途 & Exit\\
\midrule
\texttt{racergt design} & 鎖定 protocol、排程、chunk manifest & 0\\
\texttt{racergt audit} & 比對批次與鎖定 protocol & 失敗為 2\\
\texttt{racergt run} & 執行完整 pipeline & FAIL 為 3\\
\texttt{racergt simulate} & 產生受控資料，可選擇直接執行 pipeline & 0\\
\texttt{racergt export-stata} & CSV 轉 Stata 118 & 0\\
\bottomrule
\end{tabularx}
\end{table}

\end{document}
